% !TeX program = lualatex
% !TeX root    = UCIMED-Presentacion-plantilla.tex
%%%%%%%%%%%%%%%%%%%%%%%%%%%%%%%%%%%%%%%%%%%%%%%%%%%%%%%%%%%%%%%%%%%%%%%%%%%%%%%
% Plantilla de ejemplo — tema de conferencia beamerthemeUCIMEDpres.sty
% Ejercita TODAS las macros de la API (ver el bloque «PLANTILLAS DE
% REFERENCIA» al final del archivo).
%
% Compilar:  latexmk -lualatex UCIMED-Presentacion-plantilla.tex
% *** SIEMPRE con latexmk (o 2 pasadas de lualatex seguidas), NUNCA una sola
% pasada suelta *** — la portada usa tikz remember-picture/overlay, que
% exige que el .aux ya tenga el tamaño de página (limitación estándar de
% pgf/TikZ, no un defecto: beamerthemeCICANUM.sty, el original que este
% tema replica, tiene la misma exigencia). Ver la cabecera de
% beamerthemeUCIMEDpres.sty para el detalle verificado.
%
% Opciones del tema:
%   [nologos] portada e interiores sin el logo institucional
%   [nofondo] diapositivas internas sin la banda superior suave
%%%%%%%%%%%%%%%%%%%%%%%%%%%%%%%%%%%%%%%%%%%%%%%%%%%%%%%%%%%%%%%%%%%%%%%%%%%%%%%
\documentclass[aspectratio=169,10pt,spanish]{beamer}

\usepackage[spanish,es-noquoting,es-nodecimaldot]{babel}
\usepackage{csquotes}
\usepackage{fontspec}
\usepackage{amsmath,amssymb}
\usepackage{siunitx}
\usepackage{booktabs}

% Tipografía: palo seco
\IfFontExistsTF{TeX Gyre Heros}{\setsansfont{TeX Gyre Heros}}{}
\renewcommand{\familydefault}{\sfdefault}

\usepackage{colors-UCIMED}
\usepackage{comands-UCIMED}
\setclass{CBI-07} % Física II — fija \theuni, \theunishort, \theschool, \ulogo

\usepackage{beamerthemeUCIMEDpres}

\sisetup{separate-uncertainty=true,range-units=single}

%==< Metadatos =================================================================
\title{Electricidad y magnetismo en el cuerpo humano:}
\subtitle{de la ley de Ampère a la estimulación magnética transcraneal.}
\author{Prof. Luis Fernando Umaña Castro}
\institute{Escuela Autónoma de Ciencias Médicas de Centro América, UCIMED}
\date{2026}
\edicionUCIMED{\bfseries II ciclo 2026\\\bfseries Grupo 01}
% \pieUCIMED{\bfseries Escuela Autónoma de Ciencias Médicas de Centro América}

%%%%%%%%%%%%%%%%%%%%%%%%%%%%%%%%%%%%%%%%%%%%%%%%%%%%%%%%%%%%%%%%%%%%%%%%%%%%%%%
\begin{document}
%%%%%%%%%%%%%%%%%%%%%%%%%%%%%%%%%%%%%%%%%%%%%%%%%%%%%%%%%%%%%%%%%%%%%%%%%%%%%%%

%==< Portada >====================================================================
\portada

%==< Diapositiva de sección >======================================================
\seccion{Campo magnético e inducción}{De Biot-Savart a la ley de Faraday, con
  aplicaciones clínicas}

%==< Tarjetas + \punto + \subsub + tabularx-ish body >=============================
\begin{frame}{¿Por qué el electromagnetismo importa en medicina?}
  \framesubtitle{Tres puertas de entrada al tema}
  \begin{columns}[T,onlytextwidth]
    \begin{column}{0.40\textwidth}
      % \figura acota el ANCHO del minipage/imagen; ucimed.png es casi
      % cuadrado (2198x2353 px, razón 1.07) — un ancho discreto (18mm) evita
      % que domine la diapositiva (a \linewidth de esta columna, ~58mm,
      % la imagen saldría ~62mm de alto y desbordaría el frame).
      \figura[18mm]{ucimed.png}%
        {Identidad institucional UCIMED (relleno de ejemplo, sin dependencias
         externas).}
    \end{column}
    \begin{column}{0.55\textwidth}
      \punto{El reto}{explicar por qué un campo magnético variable induce
        corrientes en tejido excitable.}
      \punto{La solución}{la ley de Faraday: \(\varepsilon = -\dfrac{d\Phi_B}{dt}\).}
      \punto{La aplicación}{la estimulación magnética transcraneal (TMS) usa
        exactamente ese principio.}
    \end{column}
  \end{columns}
  \vfill
  \subsub{El principio de detección}
  \pasos{Bobina}{Pulso de corriente}{Campo \(B\) variable}%
        {Corriente inducida}{Despolarización neuronal}
\end{frame}

%==< Cadena de tres tarjetas + flecha ancha >=======================================
\begin{frame}{Del laboratorio a la clínica}
  \framesubtitle{Cadena conceptual de tres pasos}
  \vfill
  \cadenatres{\iconflask\ Ley de Ampère}{%
      Un solenoide produce un campo \(B\) uniforme y controlable dentro de
      la bobina.}%
    {\icongear\ Bobina TMS}{%
      La bobina en forma de ocho concentra el campo sobre la corteza
      motora del paciente.}%
    {\iconbolt\ Respuesta neuronal}{%
      La corriente inducida despolariza axones y produce una respuesta
      motora medible.}
  \vfill
  \flecha{De la ecuación de Maxwell-Faraday a la señal en el electromiógrafo}
\end{frame}

%==< Tabla + íconos sueltos >========================================================
\begin{frame}{Comparación de técnicas de neuroestimulación}
  \framesubtitle{Ejercita \texttt{tablaUCIMED}, \texttt{filaenc}, \texttt{filaalt}}
  % \UCIMEDhdrfmt se repite en cada celda del encabezado: en tabular/tabularx
  % un \color{}/\bfseries tecleado antes del primer `&` no se propaga a las
  % celdas siguientes (cada celda abre su propio grupo) — ver la nota en
  % beamerthemeUCIMEDpres.sty junto a \filaenc.
  \begin{tablaUCIMED}{L{28mm}YY}
    \filaenc Técnica &\UCIMEDhdrfmt Principio físico &\UCIMEDhdrfmt Invasividad \\ \midrule
    \filaalt TMS & Inducción electromagnética (Faraday) & No invasiva \\
    tDCS & Corriente directa transcraneal (Ohm) & No invasiva \\
    \filaalt DBS & Estimulación eléctrica directa & Invasiva (electrodo) \\
    \bottomrule
  \end{tablaUCIMED}
  \vfill
  \subsubi{\iconcheck}{Ventaja compartida}: ninguna requiere abrir el cráneo,
  salvo la estimulación cerebral profunda (DBS).
\end{frame}

%==< Cuatro tarjetas en fila (patrón de referencia) >================================
\begin{frame}{Cuatro mecanismos de acoplamiento campo-tejido}
  \begin{columns}[t,onlytextwidth]
    \begin{column}{0.235\textwidth}
      \tarjeta[0.9\linewidth]{Inducción}{Faraday: \(B\) variable induce \(\varepsilon\).}
    \end{column}
    \begin{column}{0.235\textwidth}
      \tarjeta[0.9\linewidth]{Conducción}{Corriente directa a través del tejido
        (ley de Ohm).}
    \end{column}
    \begin{column}{0.235\textwidth}
      \tarjeta[0.9\linewidth]{Capacitivo}{Acoplamiento vía membranas
        (bioimpedancia, Cole--Cole).}
    \end{column}
    \begin{column}{0.235\textwidth}
      \tarjeta[0.9\linewidth]{Piezoeléctrico}{Deformación mecánica del hueso
        (Fukada--Yasuda).}
    \end{column}
  \end{columns}
\end{frame}

%==< Cierre >==========================================================================
\seccion{Gracias}{Preguntas y comentarios}

%%%%%%%%%%%%%%%%%%%%%%%%%%%%%%%%%%%%%%%%%%%%%%%%%%%%%%%%%%%%%%%%%%%%%%%%%%%%%%%
\end{document}
%%%%%%%%%%%%%%%%%%%%%%%%%%%%%%%%%%%%%%%%%%%%%%%%%%%%%%%%%%%%%%%%%%%%%%%%%%%%%%%

%==================================================================================
% PLANTILLAS DE REFERENCIA (fuera del documento; copiar y pegar arriba)
%==================================================================================
% \begin{frame}{Título}{Subtítulo opcional}
%   Cuerpo en \textcolor{UCIMEDtexto}{azul pizarra institucional}; acentos con
%   \alert{azul} o \textcolor{accentPrimary}{dorado}.
%   \begin{itemize}\item Punto. \item Con unidades: \SI{37}{\celsius}.\end{itemize}
% \end{frame}
%
% Subsubtítulos sueltos:      \subsub{Título}  /  \subsubi{\iconcheck}{Título}
% Entrada compacta:           \punto{Etiqueta}{texto}
% Banda-flecha:                \flecha{texto}          \flecha[80mm]{texto}
% Flecha corta entre cajas:   \flechita
% Cadena de tres tarjetas:    \cadenatres{T1}{C1}{T2}{C2}{T3}{C3}
% Cadena de cinco pasos:      \pasos{a}{b}{c}{d}{e}
% Tarjeta suelta:             \tarjeta[0.9\linewidth]{Título}{Contenido}
% Figura con pie:             \figura[0.4\textwidth]{archivo}{pie}
% Íconos:  \iconcheck \icongear \iconflask \iconchip \iconbolt  (opt: [color])
%
% Tabla con encabezado del color principal (repetir \UCIMEDhdrfmt en cada
% celda del encabezado — \color/\bfseries no cruza el `&` en tabular/tabularx):
% \begin{tablaUCIMED}{L{20mm}YY}
%   \filaenc &\UCIMEDhdrfmt Col A &\UCIMEDhdrfmt Col B \\ \midrule
%   \filaalt Fila 1 & ... & ... \\
%   Fila 2 & ... & ... \\ \bottomrule
% \end{tablaUCIMED}
%
% Cuatro tarjetas en fila:
% \begin{columns}[t,onlytextwidth]
%   \begin{column}{0.235\textwidth}\tarjeta[0.9\linewidth]{C1}{...}\end{column}
%   ... x4 ...
% \end{columns}
%
% Diapositiva de sección:     \seccion{Texto grande}{Texto de apoyo}
% Identidad/logo:              \setclass{CBI-07} en el preámbulo (o CB-09/CBI-06);
%                               \theuni \theunishort \theschool quedan fijados.
% Metadatos de portada extra:  \edicionUCIMED{...}   \pieUCIMED{...}
